\section{A decidable axiom presentation and numerical checker for \(\PA\)}
\label{sec:p05-standard-pa-checker}

The c.e.-theory construction uses an arbitrary fixed domain index and
accepting histories.  For \(\PA\) we now specify a separate direct
axiom-checking convention. This section gives the numerical construction;
\cref{def:p05-selected-direct-pa-graph} supplies its arithmetic formula,
followed by the internal derivability proofs. Its numerical proofs are in
\cref{annex:p5.direct-pa-axioms,annex:p5.standard-pa-checker}.

Use exactly the arithmetic tree codes, deterministic capture-avoiding
substitution and twenty-rule annotated \(\mathsf{ND}_{\mathsf{FOL}}\) checker of
\cref{sec:arithmetic-coding-04}.  Expand numerals, infix notation,
comparisons and marked bounded quantifiers before choosing a matrix code.
Write \(v_i=N_0(i)\) and \(o=N_1()\)
for the variable and zero-term codes.  In the seven fixed arithmetic axioms,
write the two distinct displayed variable names \(x,y\) as \(x_0,x_1\).
This fixes their previously schematic names for literal serialization;
it does not turn their term instances into additional axiom leaves.
The seven chosen axiom codes are the following exact natural-number
constructor expressions:
\[
\begin{aligned}
 q_1&=N_{11}(0,N_7(N_5(N_2(v_0),o))),\\
 q_2&=N_{11}(0,N_{11}(1,
       N_{10}(N_5(N_2(v_0),N_2(v_1)),N_5(v_0,v_1)))),\\
 q_3&=N_{11}(0,N_{10}(N_7(N_5(v_0,o)),
                         N_{12}(1,N_5(v_0,N_2(v_1))))),\\
 q_4&=N_{11}(0,N_5(N_3(v_0,o),v_0)),\\
 q_5&=N_{11}(0,N_{11}(1,
       N_5(N_3(v_0,N_2(v_1)),N_2(N_3(v_0,v_1))))),\\
 q_6&=N_{11}(0,N_5(N_4(v_0,o),o)),\\
 q_7&=N_{11}(0,N_{11}(1,
       N_5(N_4(v_0,N_2(v_1)),N_3(N_4(v_0,v_1),v_0)))).
\end{aligned}
\]
They code, in order, the universal closures of
\(Sx_0\ne0\), \(Sx_0=Sx_1\to x_0=x_1\),
\(x_0\ne0\to\exists x_1(x_0=Sx_1)\), \(x_0+0=x_0\),
\(x_0+Sx_1=S(x_0+x_1)\), \(x_0\times0=0\), and
\(x_0\times Sx_1=(x_0\times x_1)+x_0\).
The sign \(\ne\) in this explanatory list is expanded as negated equality
in the codes.  The seven constants are not prime-power encodings.

\begin{hbdefinition}{def:p05-canonical-induction-code}{Canonical closure and induction-axiom codes}
For a formula code \(m\), let \(F(m)\) be the list of its distinct free
variable indices in increasing numerical order.  Obtain it by scanning
\(j<m\), retaining \(j\) exactly when it belongs to the computed
free-variable list.  Every occurring index of a valid expression is below
its tree code, so this scan includes every free variable.  Put \(F(m)=0\)
when \(m\) is not a formula code.

For a raw list \(P=[j_0,\ldots,j_{r-1}]\) and formula code \(c\), set
\[
 U(P,c)=N_{11}(j_0,N_{11}(j_1,\ldots N_{11}(j_{r-1},c)\ldots)),
 \qquad U(0,c)=c.
\]
Compute this by reversing \(P\), then wrapping once for each entry.
It is defined on arbitrary numerical inputs by the same list recursion;
uses as a formula constructor include a formula type test on \(c\).
Thus increasing indices appear from the outside inward.  The canonical
universal closure of a formula code \(c\) is \(U(F(c),c)\).

Let \(\mathsf{sub}(m,i,t)\) be the already specified substitution code
for replacing free \(x_i\) by the term coded by \(t\).  For a formula
code \(m\) and any \(i\in\N\), define
\[
\begin{aligned}
 b(m,i)&=\mathsf{sub}(m,i,o),\\
 s(m,i)&=\mathsf{sub}(m,i,N_2(v_i)),\\
 J(m,i)&=N_{10}\bigl(
    N_8(b(m,i),N_{11}(i,N_{10}(m,s(m,i)))),N_{11}(i,m)\bigr),\\
 \mathsf{IndAx}(m,i)&=U(F(m)\setminus[i],J(m,i)).
\end{aligned}
\]
Here \(F(m)\setminus[i]\) means filtering out the entry equal to \(i\),
preserving the list order.  If \(m\) is not a formula code, set
\(\mathsf{IndAx}(m,i)=0\).  The constructor permits \(x_i\) to be absent
freely from the matrix, and permits any bound-variable shadowing already
present there.  It uses the book's exact substitution, not textual replacement
of bound occurrences and not an alpha-equivalence class.
\end{hbdefinition}

\begin{hbdefinition}{def:p05-direct-pa-axioms}{Direct \(\PA\) axiom recognizer}
The numerical predicate \(\mathsf{Ax}^{\rm d}_{\PA}(a)\) holds exactly
when the following complete decision procedure returns one.
\begin{enumerate}
\item Reject unless \(a\) is a sentence code.  Accept if
      \(a=q_j\) for one of \(j=1,\ldots,7\).
\item Starting with \(a\), strip its maximal leading sequence of universal
      quantifiers, recording their indices in their outside-in order \(P\).
      The residual code is \(c\).  Each stripped matrix is a proper child;
      at most \(a\) stripping updates suffice.  After the first nonuniversal
      constructor, make the state fixed for all remaining updates.
\item Require \(c=N_{10}(d,N_{11}(i,m))\).  Require
      \(d=N_8(b,N_{11}(i',N_{10}(m',s)))\), with all displayed constructor
      arities exact.  Reject unless \(i'=i\), \(m'=m\), and \(m\)
      is a formula code.
\item Compute \(\mathsf{IndAx}(m,i)\) by the preceding definition.
      Accept exactly when it equals \(a\) literally.  Equivalently, check
      \(b=b(m,i)\), \(s=s(m,i)\), and
      \(P=F(m)\setminus[i]\), with literal equality of the listed data.
\end{enumerate}
Every unaccepted input has value zero.  The definition applies no additional
alpha-equivalence, logical-equivalence, term-instance or theoremhood closure.
An alpha-equivalent formula is accepted only if it independently has one
of the displayed literal forms.  Let \(\PA^{\rm d}\) be the set of
sentences whose codes this predicate accepts.
\end{hbdefinition}

\begin{hbtheorem}{thm:p05-direct-pa-axioms}{proved=annex:p5.direct-pa-axioms}{Exactness of the direct axiom presentation}
The functions and predicate just defined are primitive recursive.  The
accepted codes are exactly
\[
  \{q_1,\ldots,q_7\}\cup
 \{\mathsf{IndAx}(m,i):m\text{ is a formula code},\ i\in\N\}.
\]
They are codes of sentences.  For every arithmetic formula \(\sigma\),
\[
 \PA^{\rm d}\derives_{\mathsf{ND}_{\mathsf{FOL}}}\sigma
       \quad\Longleftrightarrow\quad
 \PA\derives_{\mathsf{ND}_{\mathsf{FOL}}}\sigma,
\]
where the seven fixed arithmetic axioms use the literal representatives
just declared, and the source induction scheme permits its displayed
parameter closures.  The equivalence includes any order of the free
parameters in those closures.  Any extra vacuous parameter binders, when
used as a source presentation variant, are transported explicitly as well.
\end{hbtheorem}
\Cref{annex:p5.direct-pa-axioms} proves the substitution identities,
recognition bounds and both source-presentation transports.

The source's free parameter order is a closing convention, not a new
induction principle.  The annex gives actual \(\forall E\) and
\(\forall I\) derivations in both directions, then verifies their insertion
into larger derivations.  In particular it checks the all-occurrence
freshness condition of \(\exists E\); closedness alone would not justify
that insertion.

Examples of rejected induction shapes are: a missing base conjunct; a
different step/conclusion binder; a mismatched copy of the matrix; substitution
that changes an occurrence beneath its own induction-variable binder; a
parameter prefix that is not strictly increasing; and an additional vacuous
prefix binder.
Reordering the two required parameter binders can therefore produce a
derivable sentence that fails the literal axiom test.  Such a derivation
uses the actual ND rules rather than another axiom-recognition convention.

\begin{hbexample}{def:p05-induction-shadowing-example}{Parameters and shadowing in induction recognition}
Let \(\theta\) be \(x_0+x_2=x_5\) and let the induction index be zero.
Its induction core is
\[
 \bigl(0+x_2=x_5\land
  \forall x_0((x_0+x_2=x_5)\to(Sx_0+x_2=x_5))\bigr)
        \to\forall x_0(x_0+x_2=x_5).
\]
The canonical axiom prefixes this core by \(\forall x_2\forall x_5\).
The reversed prefix \(\forall x_5\forall x_2\) is rejected by the
literal recognizer.  From the accepted axiom, instantiate first with
\(x_2\), then with \(x_5\), to obtain the core; generalize first
\(x_2\), then \(x_5\), to derive exactly that reversed closure.
The sole open axiom is a sentence throughout these four inferences.

For the matrix \(\theta=\forall x_0(x_0=x_2)\) and induction index
zero, both computed substitution instances equal \(\theta\) literally:
substitution stops at its outer binder.  The canonical parameter prefix
still binds \(x_2\), and the induction constructor is valid despite
the absence of a free \(x_0\).  A proposed step that replaces the bound
\(x_0\) inside this matrix by \(Sx_0\) fails the exact comparison.
Thus neither a free-occurrence requirement nor replacement of bound
occurrences belongs to the schema being recognized.
\end{hbexample}

\begin{hbdefinition}{def:p05-standard-pa-numerical-proof}{Standard direct numerical \(\PA\) proof codes}
A code in this new proof sort has the form \(N_{19}(d,0)\), where
\(d=N_{15}(R)\).  Tag \(19\) is a new outer record tag; it is distinct
from the c.e.-certificate outer tag \(16\).  The field zero is the empty
certificate field.  It does not stand for an unverified accepting history.

Define \(\mathsf{StdProof}_{\PA}(p,x)\) by this exact numerical procedure.
\begin{enumerate}
\item Require outer tag \(19\) with exactly two arguments, the second zero;
      require the first argument to have tag \(15\) with exactly one argument
      \(R\), and require \(n=\ell(R)>0\).
\item Apply the complete twenty-rule derivation procedure of
      \cref{sec:arithmetic-coding-04} to \(R\), in increasing row order.
      Each row has exactly \([r,a,J,U,l]\); rule tags are precisely
      \(0,\ldots,19\), with their exact premise arities, parameter lists,
      formula/term tests, nonassumption marker zero, and earlier references.
      Compute each open-occurrence list \(O_j\) by the stated branch-specific,
      formula-specific, nonvacuous discharge filters.  Apply the stated free-
      versus all-variable tests and the exact substitution/alpha tests.
\item Require every nonlast row to have exactly one parent occurrence and
      the last row none.  Require a sentence at the last row and literal
      equality of its code with \(x\).
\item For every leaf index in \(O_{n-1}\), apply
      \(\mathsf{Ax}^{\rm d}_{\PA}\) to its row formula and require value one.
      Temporary assumptions absent from this final open list need no axiom
      test.  An empty open list passes this step.
\end{enumerate}
The relation holds exactly when all these checks pass; its characteristic
function returns one in that case and zero otherwise.
There is no Comp call and no implicit replacement of a noncanonical axiom
by a theorem at a leaf.  All twenty inference-rule tests themselves are
unchanged, including body-order implication elimination and premise-free
identity introduction.
\end{hbdefinition}

\begin{hbexample}{def:p05-direct-pa-proof-examples}{Axiom and empty-context direct proofs}
For each \(q_j\), the one-row derivation record
\(d_j=N_{15}([[0,q_j,0,0,0]])\) is an assumption with one open leaf.
The direct code \(N_{19}(d_j,0)\) proves that arithmetic axiom: the leaf
passes the fixed-constant axiom test, and no computation certificate is needed.

Put \(e=N_5(o,o)\).  The one-row record
\(d_{=}=N_{15}([[18,e,0,[o],0]])\) is the premise-free \(=I\) proof
of \(0=0\).  Its open list is empty, and \(N_{19}(d_{=},0)\) passes
the direct checker with root code \(e\).  Replacing the empty outer
field in either example by a nonempty list makes the direct code fail.
The empty-context example does not test membership of \(0=0\) in the
nonlogical axiom set.
\end{hbexample}

\begin{hbtheorem}{thm:p05-standard-pa-numerical-proof}{proved=annex:p5.standard-pa-checker}{Numerical standard \(\PA\) proof meaning}
The relation \(\mathsf{StdProof}_{\PA}\subseteq\N^2\) is primitive
recursive, and for every arithmetic sentence \(\sigma\),
\[
 \PA\derives\sigma
   \quad\Longleftrightarrow\quad
 \exists p\in\N\,\mathsf{StdProof}_{\PA}
                         (p,\ulcorner\sigma\urcorner).
\]
For any fixed c.e. domain-index convention for the source \(\PA\) axioms,
the right-hand side is also equivalent to existence of a certified
\(\mathsf{Proof}_{\PA}\) code from
\cref{def:fixed-proof-predicate}.  This is an equivalence of proof existence;
the two predicates and their code sorts are distinct.
\end{hbtheorem}
\Cref{annex:p5.standard-pa-checker} proves axiom recognition, source-presentation
transport, the primitive-recursion bounds, and both proof-meaning directions.
It also specifies the bridge to the separate c.e.-certificate convention.

\paragraph{The arithmetic formula and its internal verification.}
The numerical checker supplies the procedure whose fixed description is
expanded in \cref{def:p05-selected-direct-pa-graph}, using the graph
construction of \cref{def:p05-pa-pr-graph}. Its exact formula is
\(\mathsf{Prf}^{19}_{\PA}\), with existential closure
\(\mathsf{Prv}^{19}_{\PA}\).
The internal proof-combining theorem is
\cref{thm:p05-direct-pa-internal-mp}; the full original-sentence
evidence construction is \cref{thm:p05-original-sigma-evidence}, and
\cref{thm:p05-direct-pa-d3} proves positive introspection for that
same predicate. These later verifications supply the internal results
in addition to the numerical proof-meaning theorem above.
